%%%%%%%%%%%%%%%%%%%%%%%%%%%%%%%%%%%%%%%%%%%%%%%%%%%%%%%%%
%Modelo de como colocar formulas directamente sobre el texto, o bien Entre párrafos.

%Para mas información sobre las figura comprobar la documentación del paquete "amsmath".
%%%%%%%%%%%%%%%%%%%%%%%%%%%%%%%%%%%%%%%%%%%%%%%%%%%%%%%%%

Para escribir ecuaciones en medio de un texto tenemos que encerrar la ecuación entre símbolos de dolar $E=mc^2$. También podemos hacer colocar las formulas en huecos centrado entre párrafos de las dos siguientes formas:

\[E=mc^2\]

\begin{equation}
    \label{energy}
    E=mc^2
\end{equation}

Con el segundo formato vemos que podemos numerar las ecuaciones cómodamente y referenciarlas en el texto de forma simple, con le comando \ref{energy}.

Tal como hemos visto, para crear superíndices usamos $a^b$, mientras que para crear subíndices usamos $a_b$. Esto lo podemos llevar al extremo creando formular realmente complejas, como la siguiente:

\[ T^{i_1 i_2 \dots i_p}_{j_1 j_2 \dots j_q} = T(x^{i_1},\dots,x^{i_p},e_{j_1},\dots,e_{j_q}) \]

Para escribir fracciones usamos $\frac{a}{b}$, mientras que para integrales usamos $\int_0^1$, combinándolo con los superíndices y subíndices. Con esto podemos crear formulas complejas como la siguiente:

\[ \int_0^1 \frac{1}{e^x} =  \frac{e-1}{e} \]

Las letras griegas se pueden escribir en minúscula ($\omega$, $\delta$) o en mayúscula ($\Omega$, $\Delta$).

Los operadores especiales, tales como senos, consenos o logaritmos se escriben con los comandos ($\sin(\beta)$, $\cos(\alpha)$, $\log(x)$).